% ============================================================
% LOGICAL ENTITY IDENTIFIERS
% Viability Philosophy System
% Book I: Ontology of Viability
% ============================================================

\cleardoublepage

\section*{Logical Entity Identifiers}

\phantomsection
\addcontentsline{toc}{section}{Logical Entity Identifiers}

% ------------------------------------------------------------
% 1. PURPOSE AND SCOPE
% ------------------------------------------------------------

\subsection*{1. Purpose and Scope}

Logical entity identifiers provide a system for
identifying conceptual, theoretical, and
argumentative entities within the Viability
Philosophy System.

Their purpose is to establish persistent identities
for entities that may be introduced, defined,
examined, developed, qualified, or applied across
different textual locations and, where applicable,
across different books of the system.

A logical entity identifier identifies an entity
as such. It does not identify the chapter, section,
subsection, paragraph, page, or source file in which
the entity happens to be discussed.

This distinction separates the identity of an
entity from the location of its presentation.

An entity may appear in several locations without
acquiring a new logical identity each time it is
mentioned. Its identifier remains associated with
the entity, while each textual occurrence is
located through the structural referencing system.

Logical entity identifiers therefore serve as
persistent references to the intellectual
architecture of the system.

They support:

\begin{itemize}
    \item the unambiguous identification of logical
    entities;
    \item the distinction between entities that
    may otherwise be referred to by similar names;
    \item references to entities across chapters
    and books;
    \item the tracing of conceptual and argumentative
    dependencies;
    \item the coordination of definitions, propositions,
    theories, hypotheses, and other logical entities;
    \item the maintenance of conceptual continuity
    during manuscript revision.
\end{itemize}

% ------------------------------------------------------------
% 2. GENERAL FORMAT
% ------------------------------------------------------------

\subsection*{2. General Format}

A logical entity identifier consists of a
category prefix followed by a numerical sequence.

The general format is:

\begin{center}
    \texttt{XX-001}
\end{center}

Here, \texttt{XX} represents the prefix assigned
to a logical entity family, while \texttt{001}
represents the entity's numerical identifier
within that family.

Examples of the established identifier format
include:

\begin{center}
\begin{tabular}{ll}
    \toprule
    \textbf{Identifier} & \textbf{Identifier Family} \\
    \midrule
    \texttt{AX-001} & AX \\
    \texttt{DF-001} & DF \\
    \texttt{PR-001} & PR \\
    \texttt{TH-001} & TH \\
    \texttt{HY-001} & HY \\
    \texttt{AR-001} & AR \\
    \texttt{PF-001} & PF \\
    \bottomrule
\end{tabular}
\end{center}

The prefix identifies the family to which
the entity belongs. The numerical component
distinguishes the individual entity within
that family.

The complete identifier must be treated
as a single identification unit.

% ------------------------------------------------------------
% 3. ESTABLISHED IDENTIFIER FAMILIES
% ------------------------------------------------------------

\subsection*{3. Established Identifier Families}

The logical entity identification system
currently recognizes seven identifier
families:

\begin{center}
\begin{tabular}{ll}
    \toprule
    \textbf{Prefix} & \textbf{Example} \\
    \midrule
    \texttt{AX} & \texttt{AX-001} \\
    \texttt{DF} & \texttt{DF-001} \\
    \texttt{PR} & \texttt{PR-001} \\
    \texttt{TH} & \texttt{TH-001} \\
    \texttt{HY} & \texttt{HY-001} \\
    \texttt{AR} & \texttt{AR-001} \\
    \texttt{PF} & \texttt{PF-001} \\
    \bottomrule
\end{tabular}
\end{center}

These prefixes constitute distinct
identifier families within the system.

Their classification must be preserved
consistently throughout the manuscript
and across the books of the Viability
Philosophy System.

The precise semantic definition, scope,
and admissibility conditions of each
family must be governed by the formal
classification adopted for the system.

A prefix must not be assigned to an
entity merely on the basis of an informal
similarity of wording or an isolated
textual occurrence.

The assignment of an entity to a family
requires consistency with the conceptual
and logical classification of the system.

% ------------------------------------------------------------
% 4. GLOBAL UNIQUENESS
% ------------------------------------------------------------

\subsection*{4. Global Uniqueness}

Logical entity identifiers are globally
unique within the Viability Philosophy
System.

Uniqueness is not restricted to a single
chapter, part, manuscript file, or book.

Once an identifier has been assigned
to a logical entity, the same identifier
must not be assigned to an unrelated
entity elsewhere in the system.

For example, if \texttt{AX-001} is
assigned to a particular entity in
Book I, another entity in Book II
must not independently receive
the identifier \texttt{AX-001}.

The identifier belongs to the entity
within the system as a whole, not
merely to the book in which the entity
is first introduced.

This principle ensures that logical
references remain unambiguous across
the entire system.

% ------------------------------------------------------------
% 5. SEPARATION FROM STRUCTURAL
% IDENTIFIERS
% ------------------------------------------------------------

\subsection*{5. Separation from Structural
Identifiers}

Logical entity identifiers must be
distinguished from structural identifiers.

A structural identifier specifies
the location of a textual element
within the hierarchy of a book.

A logical entity identifier specifies
the identity of an entity discussed
or established in the text.

For example:

\begin{center}
\begin{tabular}{ll}
    \toprule
    \textbf{Identifier} & \textbf{Function} \\
    \midrule
    \texttt{VS.B01.P01.C01.S01.SS01.P001}
        & Structural location \\
    \texttt{AX-001}
        & Logical entity identity \\
    \bottomrule
\end{tabular}
\end{center}

The structural identifier may change
when the corresponding textual element
is relocated within the manuscript.

The logical entity identifier, however,
must remain associated with the same
entity independently of such relocation.

The two systems are complementary.
Neither can replace the other.

% ------------------------------------------------------------
% 6. IDENTITY AND TEXTUAL OCCURRENCE
% ------------------------------------------------------------

\subsection*{6. Identity and Textual Occurrence}

A logical entity may occur in multiple
places throughout the book.

Its initial introduction, formal
definition, subsequent analysis,
application, criticism, or revision
may take place in different textual
locations.

These occurrences do not automatically
constitute different logical entities.

A single entity retains its logical
identifier across its relevant
occurrences, provided that its identity
has not been replaced by a genuinely
distinct entity.

For this reason, the identifier must
not be assigned anew whenever an
existing entity is mentioned.

The identification system distinguishes
between:

\begin{itemize}
    \item the logical entity itself;
    \item the textual occurrence in which
    the entity is discussed;
    \item the structural location of
    that occurrence.
\end{itemize}

These are related but non-identical
aspects of the referencing system.

% ------------------------------------------------------------
% 7. ASSIGNMENT OF IDENTIFIERS
% ------------------------------------------------------------

\subsection*{7. Assignment of Identifiers}

A logical entity identifier should
be assigned when an entity is admitted
into the formal logical inventory
of the Viability Philosophy System.

The assignment must follow the
established family classification
and uniqueness rules.

The following principles govern
the process.

\subsubsection*{7.1. Classification Before Assignment}

The appropriate identifier family
must be determined before an identifier
is assigned.

An entity must not be placed in a
family solely because a particular
prefix appears convenient.

Classification must follow the
formal conceptual and logical
distinctions of the system.

\subsubsection*{7.2. Numerical Assignment}

The numerical component must
distinguish the entity from all
other entities within the same
identifier family.

Numbers must be assigned consistently
and must not be duplicated.

The established three-digit format
should be preserved:

\begin{center}
    \texttt{001}, \texttt{002},
    \texttt{003}, \ldots
\end{center}

Leading zeros are part of the
identifier format and must not
be omitted.

\subsubsection*{7.3. No Assignment by Mere Mention}

A term, phrase, or conceptual expression
does not automatically become a
separately identified logical entity
merely because it appears in the text.

An identifier should be assigned
only when the expression refers
to an entity that is recognized
as such within the logical
organization of the system.

\subsubsection*{7.4. No Duplicate Assignment}

Before assigning a new identifier,
the existing logical entity inventory
must be checked.

If the intended entity has already
been assigned an identifier, the
existing identifier must be reused
rather than creating a duplicate.

% ------------------------------------------------------------
% 8. STABILITY AND REVISION
% ------------------------------------------------------------

\subsection*{8. Stability and Revision}

Logical entity identifiers are
intended to preserve identity
throughout the development of
the system.

Changes in wording, exposition,
structural location, or the order
of presentation do not, by themselves,
justify the creation of a new
logical entity identifier.

The relevant question is whether
the entity itself remains the
same entity within the conceptual
and logical organization of
the system.

\subsubsection*{8.1. Editorial Revision}

An editorial revision that leaves
the identity of an entity intact
does not require a new identifier.

The existing identifier must
be retained.

\subsubsection*{8.2. Structural Relocation}

Moving an entity's discussion
from one chapter or section to
another does not alter its
logical identity.

Its structural location references
may need to be updated, but
its logical entity identifier
remains unchanged.

\subsubsection*{8.3. Conceptual Modification}

A substantial conceptual modification
requires examination of whether
the entity retains its identity.

A modification must not automatically
be treated as either a continuation
or a replacement.

The decision must be made according
to the actual conceptual and logical
relationship between the earlier
and revised formulations.

\subsubsection*{8.4. Replacement and Retirement}

If an entity is replaced by a
distinct entity, the distinction
must be reflected in the logical
entity inventory.

An identifier that has already
been assigned must not be
reassigned to an unrelated entity.

Where an entity is no longer
used, its identifier should
remain traceable in the historical
record of the system.

% ------------------------------------------------------------
% 9. RELATION TO LOGICAL DEPENDENCIES
% ------------------------------------------------------------

\subsection*{9. Relation to Logical Dependencies}

Logical entity identifiers provide
the reference points through which
relationships and dependencies
between entities can be specified.

A logical dependency must refer
to the identifiers of the entities
involved rather than relying
exclusively on their names or
the locations where they are
discussed.

This makes it possible to distinguish
the identity of an entity from
its relations to other entities.

The identification of an entity
does not, by itself, establish
the nature or direction of its
dependencies.

Those relationships must be
specified separately through
the logical dependency system.

% ------------------------------------------------------------
% 10. RELATION TO INTERNAL
% CROSS-REFERENCES
% ------------------------------------------------------------

\subsection*{10. Relation to Internal
Cross-References}

Internal cross-references may
refer either to structural locations
or to logical entities.

These two forms of reference
must not be conflated.

A reference to a chapter, section,
subsection, or paragraph identifies
a textual location.

A reference to a logical entity
identifies the entity independently
of the particular passage in which
it is presented.

Where both forms of reference
are necessary, they may be used
together.

The conventions for constructing
and maintaining internal
cross-references are established
in the corresponding referencing
document.

% ------------------------------------------------------------
% 11. RELATION TO EXTERNAL SOURCES
% ------------------------------------------------------------

\subsection*{11. Relation to External Sources}

Logical entity identifiers are
reserved for entities belonging
to the internal logical organization
of the Viability Philosophy System.

External authors, books, articles,
documents, and other bibliographic
sources do not receive logical
entity identifiers merely because
they are cited or discussed.

External sources retain their
own bibliographic identities.

Their references must be managed
through the external source and
bibliographic referencing systems.

A logical entity in the Viability
Philosophy System may be discussed
in relation to an external source,
but the source itself must not
be confused with the internal
entity.

% ------------------------------------------------------------
% 12. MAINTENANCE OF THE LOGICAL
% ENTITY INVENTORY
% ------------------------------------------------------------

\subsection*{12. Maintenance of the Logical
Entity Inventory}

The logical entity identification
system requires a consistent
inventory of assigned identifiers.

The inventory should make it
possible to determine:

\begin{itemize}
    \item which identifiers have
    been assigned;
    \item which family each identifier
    belongs to;
    \item which logical entity
    each identifier designates;
    \item whether an identifier
    is active, revised, superseded,
    or retired;
    \item where the entity is
    introduced and discussed;
    \item which other entities
    are logically related to it.
\end{itemize}

The inventory must be maintained
in coordination with the manuscript
and its referencing documents.

An identifier must not be treated
as available merely because its
entity is absent from a particular
chapter or source file.

The inventory concerns the system
as a whole.

% ------------------------------------------------------------
% 13. SUMMARY
% ------------------------------------------------------------

\subsection*{13. Summary}

Logical entity identifiers establish
persistent identities for entities
within the Viability Philosophy
System.

Their general format is:

\begin{center}
    \texttt{XX-001}
\end{center}

The established identifier families
are:

\begin{center}
\begin{tabular}{ll}
    \toprule
    \textbf{Family} & \textbf{Example} \\
    \midrule
    \texttt{AX} & \texttt{AX-001} \\
    \texttt{DF} & \texttt{DF-001} \\
    \texttt{PR} & \texttt{PR-001} \\
    \texttt{TH} & \texttt{TH-001} \\
    \texttt{HY} & \texttt{HY-001} \\
    \texttt{AR} & \texttt{AR-001} \\
    \texttt{PF} & \texttt{PF-001} \\
    \bottomrule
\end{tabular}
\end{center}

The system is governed by the
principles of global uniqueness,
classification consistency,
identity preservation, and
separation from structural
identification.

Logical entity identifiers
remain associated with their
entities across textual locations
and books.

They constitute the identity
layer of the internal referencing
system and must be coordinated
with structural identifiers,
internal cross-references,
entity location references,
logical dependencies, and
bibliographic organization.

\cleardoublepage

% ============================================================
% END OF LOGICAL ENTITY IDENTIFIERS
% ============================================================